\scanpage{321}
\mspage{237}
% archive sheet number: 320

\uline{4.9.6.9}. Démonstration de 4.9.6.3, \uline{partie 2\textsuperscript{o}}.

C'est essentiellement ``la $=$'' démonstration que pour la partie 1\textsuperscript{o}).

$a_{\unsure{\pi|\kappa}})\Rightarrow b_\pi)$. On peut court-circuiter la démonstration directe, en regardant, pour un ensemble $A$ de cardinal $\leq\pi$, le diagramme

\diagram{\begin{tikzcd}[ampersand replacement=\&]
I\arrow[r,"\delta_A"]\arrow[d,"f"']\&I^A\arrow[d,"f^A"]\\
I'\arrow[r,"\delta_{A'}"']\&I'^{A}.
\end{tikzcd}}

Comme $f$ est cofinal, \deleted{\unsure{$f$}} $I$ filtrante, $f^A$ est aussi cofinal (4.9.\unsure{8|6}.5), d'autre part $I$ \unsure{grde|gde} devant $\pi$ signifie que $\delta_A$ est cofinal ($\forall\, A$ avec $\operatorname{card} A\leq\pi$). Donc $f^A\delta_A=\delta_{A'}f$ est cofinal, donc $\delta_{A'}f$ et $f$ le sont, donc $\delta_{A'}$ aussi, ce qui signifie que $I'$ est grde devant $\pi$.

Reste à prouver que pour $i'$ dans $I'$, $ {}_{i'}\backslash I$ est grde devant $\pi$. Cela résulte du fait que $ {}_{i'}\backslash I\xrightarrow{\unsure{q|\varphi}}I$ est coétale, \inserted{que $ {}_{i'}\backslash I$ est $0$-connexe,} et du

\uline{Lemme} 4.9.6.9.1. Soit $I$ \unsure{une} catégorie filtrante (\inserted{resp.} grde devant $\pi$). \deleted{\unsure{et}} Alors pour \unsure{tte|tout} $J$ sur $I$ coétale sur $I$, \inserted{et $0$-connexe,} $J$ est filtrante (resp. grde devant $\pi$).

\deleted{Il faut vérifier} \deleted{\unsure{qui résulte}}

Le cas filtrant est déjà connu (cf II, \ldots). Pour le cas ``grd devant $\pi$'', on va utiliser le

\uline{\deleted{\unsure{Lemme}}Corollaire} 4.9.6.9.2. Soit $J$ une catégorie. Pour que $J$ soit grde devant $\pi$, il

